% Created 2019-03-19 Tue 09:20
\documentclass[11pt]{article}
\usepackage[utf8]{inputenc}
\usepackage[T1]{fontenc}
\usepackage{fixltx2e}
\usepackage{graphicx}
\usepackage{longtable}
\usepackage{float}
\usepackage{wrapfig}
\usepackage{rotating}
\usepackage[normalem]{ulem}
\usepackage{amsmath}
\usepackage{textcomp}
\usepackage{marvosym}
\usepackage{wasysym}
\usepackage{amssymb}
\usepackage{hyperref}
\tolerance=1000
\usepackage{csquotes}
\usepackage{graphicx}
\usepackage{amsmath}
\usepackage{mathtools}
\usepackage{algpseudocode}
\usepackage[margin=0.75in]{geometry}
\usepackage{nth}
\usepackage{tikz}
\usetikzlibrary{positioning,chains,fit,shapes,shapes.misc,calc,arrows,patterns,decorations.pathreplacing}
\usepackage{ccicons}
\usepackage{calc}
\usepackage{fancyhdr}
\usepackage{fancyvrb}
\usepackage{hyperref}
\usepackage{nopageno}
\usepackage{enumitem}
\usepackage{multicol}
\usepackage{braket}
\algnewcommand\algorithmicforeach{\textbf{for each}}
\algdef{S}[FOR]{ForEach}[1]{\algorithmicforeach\ #1\ \algorithmicdo}
\setlength{\columnseprule}{1pt}
\newcommand{\bonusmarksamt}[1]{\textbf{[#1~BONUS~marks]}}
\newcommand{\bonusmarkamt}[0]{\textbf{[1~BONUS~mark]}}
\newcommand{\marksamt}[1]{\textbf{[#1~marks]}}
\newcommand{\markamt}[0]{\textbf{[1~mark]}}
\newcommand{\fillinMC}[1]{\fillinMCmath{\mbox{#1}}}
\newcommand{\fillinMCmath}[1]{\begin{tikzpicture}\draw circle [radius=0.5em];\end{tikzpicture}\ #1}
\newcommand{\fillinMCsoln}[1]{\fillinMCmathsoln{\mbox{#1}}}
\newcommand{\fillinMCmathsoln}[1]{\begin{tikzpicture}\draw[black, fill=blue] circle [radius=0.5em];\end{tikzpicture}\ #1}
\newcommand{\fillinMCAll}[1]{\fillinMCAllmath{\mbox{#1}}}
\newcommand{\fillinMCAllmath}[1]{\begin{tikzpicture}\draw (0,0) rectangle (1em,1em);\end{tikzpicture}\ #1}
\newcommand{\fillinMCAllsoln}[1]{\fillinMCAllmathsoln{\mbox{#1}}}
\newcommand{\fillinMCAllmathsoln}[1]{\begin{tikzpicture}\draw[black, fill=blue] (0,0) rectangle (1em,1em);\end{tikzpicture}\ #1}
\newcommand{\fillinblank}[1]{\fillinblankmath{\mbox{#1}}}
\newcommand{\fillinblankmath}[1]{\begingroup\setlength{\fboxsep}{1em}\setlength{\fboxrule}{2pt}\fbox{\LARGE\phantom{#1}}\endgroup}
\newcommand{\fillinblanksoln}[1]{\fillinblankmathsoln{\mbox{#1}}}
\newcommand{\fillinblankmathsoln}[1]{\begingroup\setlength{\fboxsep}{1em}\setlength{\fboxrule}{2pt}\fbox{{#1}}\endgroup}
\newcommand{\answerbox}[1][3\baselineskip]{%
\noindent\setlength{\fboxrule}{2pt}\framebox[0.7\linewidth]{%
\raisebox{2pt}[0pt][#1]{}%
}\par\medskip%
}
\usepackage{collcell,array}
\newcommand{\shiftdown}[1]{\smash{\raisebox{-.5\normalbaselineskip}{#1}}}
\newcolumntype{C}{>{\collectcell\shiftdown}c<{\endcollectcell}}
\newcommand{\shiftup}[1]{\smash{\raisebox{.5\normalbaselineskip}{#1}}}
%%% Adding Colour to Questions and Answers                                                            
\usepackage{color}
%Questions                                                                                            
\def\question#1{\red{#1}}
% Answers                                                                                             
%Answer in short form                                                                                 
\def\ans#1{\par\blu{Answer: #1}}%With header of Answer                                                
\def\soln#1{\par\blu{#1}} %Without header of Answer                                                   
% Colors                                                                                              
\definecolor{blu}{rgb}{0,0,0.5}
\def\blu#1{{\color{blu}#1}}
\definecolor{gre}{rgb}{0,.3,0}
\def\gre#1{{\color{gre}#1}}
\definecolor{red}{rgb}{0.5,0.0,0}
\def\red#1{{\color{red}#1}}
\def\norm#1{\|#1\|}
%%% End for Colours            
\def\sequence#1#2{\dsequence{#1}{1}{#2}}
\def\dsequence#1#2#3{#1_{#2}, \ldots, #1_{#3}}
\newtheorem{fact}{Fact}
\def\imax{$\texttt{S}_{\max}$}
\def\iprevmax{previous\_$\texttt{S}_{\max}$}
\def\deleted{\texttt{covered}}
%\fancyhead{}\fancyfoot[C]{\tiny This work is licensed under a \href{http://creativecommons.org/licenses/by/4.0/}{Creative Commons Attribution 4.0 International License. \ccby}\\ For license purposes, the author is the University of British Columbia.}\pagestyle{fancyplain}
\makeatletter
\def\@myproof#1[#2]{\trivlist \item [\hskip\labelsep{\bf #1 {#2}: }]}
\def\@begproof#1{\@ifnextchar[{\@myproof{#1}}{\@myproof{#1}[\unskip$\,$]}}
\newenvironment{solution}{\vskip -10pt\leavevmode\@begproof{Solution}}{\endtrivlist}
\def\sopt{S_{\textit{\scriptsize opt}}}
\makeatother
\fancyhead{}\fancyfoot[C]{\tiny This work is licensed under a \href{http://creativecommons.org/licenses/by-nc/4.0/}{Creative Commons Attribution-NonCommercial 4.0 International License. \ccbync}\\ For license purposes, the author is the University of British Columbia.}\pagestyle{fancyplain}
\usepackage{mdframed}
\BeforeBeginEnvironment{minted}{\begin{mdframed}}
\AfterEndEnvironment{minted}{\end{mdframed}}
\date{}
\title{CPSC 320 2018W2: Assignment 5 (Last assignment!)}
\hypersetup{
  pdfkeywords={},
  pdfsubject={},
  pdfcreator={Emacs 25.2.1 (Org mode 8.2.10)}}
\begin{document}

\maketitle
\vspace{-.7in}

Please follow the guidelines given in Assignment 1 for submission to
Gradescope, and for group collaboration. Remember to provide short
justifications for your answers. Submit by the deadline \textbf{Monday 
April 1, 2019 at 10PM}.

Use \LaTeX to prepare your answers. Easiest will be to use the .tex
file provided.  Please enclose each paragraph of your solution in
\verb~\soln{Your solution here...}~.  \soln{Your solution will then appear
in dark blue}, making it a lot easier for TAs to find the parts that
you wrote.

\section{Subset sums mod $n$}
\label{sec-1}
You've now seen examples of dynamic programming problems where the
recurrence describes a number (e.g., the length of a longest common
subsequence), and one example where the recurrence value is Boolean (true or false, indicating
whether or not a there is a subset of a set items whose sum is a given value).
Here you'll work with a problem where the recurrence describes a set.

Let $S = \{x_1, x_2, \ldots, x_n\}$ be a set of nonnegative integers.
We say that a value $v$, with $0 \le v \le n-1$, is \emph{feasible} with respect
to $S$ if for some non-empty subset $R$ of $S$, 
\[ 
\sum_{x \in R} x = v \pmod n.  
\]

\begin{enumerate}
\item For $1 \le i \le n$, let $V(i)$ be the set of feasible values that can
be obtained from $\{x_1, x_2, \ldots, x_i\}$.
\question{Explain why the following recurrence holds for $i \ge 1$, if we define
$V(0) = \emptyset$ (the empty set):}
 \[
    V(i) = V(i-1) + \cup_{v \in V(i-1)} \{ v + x_i \pmod n \} \cup \{ x_i \pmod n \}.
    \]
\soln{
   We can see this recurrence in three parts, the first part $V(i-1)$ ensure that $V(i)$ contains
   the $v_{\text{s}}$ for all subsets of the previous set $\{x_1, x_2 ... x_{i-1}\}$. \\ \\
   With $\cup_{v \in V(i-1)} \{ v + x_i \pmod n \}$ we are addding all the $v_{\text{s}}$ for all the 
   combinations of the previous subsets in $V(i-1)$ with our new value $x_i$, since we only care 
   about the $\pmod n$ of the sum of each subset, we dont need to re calculate the sum over each 
   subset to include $x_i$ we just need the corresponding $v$ of each subset then if we add $x_i$ to 
   it and compute $\pmod n$ it would yield the same result without having to sum over the subset $R$, 
   with this recurrence we are progressively adding to the subsets our new values $x_i$ on each iteration. \\ \\
   Lastly we are missing one subset, we have all the subsets that not include $x_i$ from the 
   first part, and all the subsets that include $x_i$ and one or more $x_j$ in $S$ with $j < i$. 
   But we have not considered the subset $R = \{x_i\}$ and that is what the last part of the 
   recurrence do, $\{ x_i \pmod n \}$, the sum of that set is just $x_i$ so we have include 
   $x_i \pmod n$ into our set $V(i)$. \\
}
\item \question{Design an algorithm that, given a set $S$ of $n$ nonnegative
integers, determines whether 0 is feasible with respect to $S$.} The
algorithm outpus "Yes" or "No".  Your algorithm should run in
$O(n^2)$ time. Your pseudocode can use set options such as $\cup$ as
in the recurrence above.
\soln{
\begin{tabbing}
xxxxxxx\=xxx\=xxx\=xxx\=xxx\=xxxx\= \kill
\> {\bf Algorithm} Subset-Sum-mod-n($S = \{x_1,x_2,\ldots,x_n\}$) \\
\>\> // Returns "Yes" if for some subset $R$ of $S$, $\sum_{x \in R} x = 0 \pmod n$   \\
\>\> Create an array Soln$[0..n]$ \\
\>\> Set Soln$[0] = \emptyset$ \\
\>\> for $i = 1$ to $n$: \\
\>\>\> Soln$[i]$ = \{S$[i] \mod n$\} \\
\>\>\> for $j = 1$ to $|$Soln$[i-1]|$: \\
\>\>\>\> $v$ = Soln$[i-1][j]$ \\
\>\>\>\> Soln$[i]$ = Soln$[i]$ $\cup$ $\{v, v+s[i] \pmod n\}$ \\
\end{tabbing}
}
\item Suppose that indeed there is a subset $R$ of $S$ such that $\sum_{x
   \in R} x = 0 \pmod n$.  \question{Write an algorithm that finds such a subset
$R$}.  You can assume that an array Soln$[0..n]$ has already been
pre-computed, where Soln$[i]$ stores the set $V(i)$ defined in part 1,
and your algorithm can use the array Soln. Your algorithm should
run in time $O(n^2)$.
\end{enumerate}
% JATJ: Complete

\section{NP True or False}
\label{sec-2}
Let $X$ and $X'$ be decision problems, where both problems
have "Yes" instances and "No" instances.
\question{State whether you think each of the following statements must be true,
must be false, or is an open question. Justify your answer.}


\begin{enumerate}
\item \textbf{Statement:}
   If $X \le_p X'$ and $X$ is \emph{not} in NP, then $X'$ is not in NP.
\soln{
   True, if $X'$ is at least as hard as $X$ which is not in NP and also not in P, then it must be that 
   $X'$ has no solution neither in P nor NP. If it was the case that there was a solution for 
   $X'$ that could be checked in polynomial then $X$ would also have a solution that could be checked 
   in polynomial time, but it is not the case since we know that $X$ is not in NP. Hence there is no 
   way to check a solution for $X'$ on polynomial time, therefore $X'$ is not in NP.
   % Open question, we can't confirm that $X'$ is not in NP by just proving that $X'$ which is not in NP is 
   % at least as hard as $X$. This can be clearer if we assign some problem to $X$, like checking if an 
   % array contain a provided element, this would be a linear problem. Knowing that $X'$ is at least as 
   % hard as finding an element in an array, don't provide us with more information about $X'$ it could 
   % easily be infinitely more difficult than P but we don't know.
}
\item \textbf{Statement:}
If $X \le_p X'$, $X$ is in P and $X'$ is in NP, then $X'$ must be in P.
\soln{
   Open question, if I can confirm or deny this, then I would win \$1 million US. By proving this 
   we would prove that P = NP.
}

\item Let
$\overline{X}$ be the \emph{complement} of problem $X$. That is,
instance $I$ of $X$ is a No-instance of $\overline{X}$ if and only
if $I$ is a Yes-instance of $X$. 
For example, the problem $\overline{\mbox{SAT}}$
is the set of \emph{unsatisfiable} 3-Sat Boolean formulas---formulas
for which there is no satisfying truth assignment.
Define $\overline{X'}$ similarly in terms of $X'$. The statement then is:

\textbf{Statement:} If $X \le_p X'$ then $\overline{X} \le_p \overline{X}'$.
\soln{
   False, even if we know that $X \le_p X'$ we can't confirm that proving a not instance is reducible 
   in the same way. So even though $\overline{X}$  has the same context as $X$ they can have complexities, 
   sometimes proving things are no-instances is faster than proving solutions are a yes-instances and 
   vice versa.
}
\end{enumerate}

\section{Common superstrings}
\label{sec-3}
A \emph{common superstring} of the strings $s_1, s_2, \ldots, s_n$ is a
string $T$ such that each $s_i$ is a substring of $T, 1 \le i \le
n$. (Note that the letters of the substring must appear consecutively in
$T$, unlike our definition of a subsequence.) For example, if the
strings are $s_1 = 101, s_2 = 0111$, and $s_3 = 010$ then the strings
010111 and 0111010 are both common superstrings of $s_1, s_2$, and
$s_3$.  

An instance of the \emph{Shortest Common Superstring} problem is a set of
strings $\{s_1, s_2, \ldots, s_n\}$ and a number $K$.
The problem is to determine whether there is a
common superstring of $s_1, s_2, \ldots, s_n$ that has length at most
$k$.  Assume that none of the strings $s_i$ is a substring of another
string $s_j$ (otherwise $s_i$ can be removed from the list of strings
without changing the length of the shortest common superstring). This
problem has been important in the context of developing algorithms for
genome sequence assembly.

An instance of the \emph{Traveling Salesperson Problem} (TSP) is a set $\{
C_1, C_2, \ldots, C_n \}$ of cities, nonnegative integer costs
$c_{ij}$ of traveling from city $C_i$ to city $C_j$ (or from $C_j$ to
$C_i$; the cost is the same either way), and a number $K$.  The
problem is to determine whether there is a tour of the cities of cost
at most $K$.  A \emph{tour} is a path that starts at some city, say $C_1$,
travels through every other city exactly once, and returns to the
starting city $C_1$. The tour cannot go through a city more than once
(except for $C_1$ which appears both at the start and at the end of
the tour).

\begin{enumerate}
\item \question{Describe a polynomial time reduction from the Shortest Common
Superstring problem to the traveling salesperson problem (TSP).}

\item \question{Is the Shortest Common Superstring problem in NP? Justify your answer.}

\item \question{Does the reduction of part 1, plus your answer to part 2, imply
that the Shortest Common Superstring problem is NP-complete?  Why
or why not?}
\end{enumerate}

\section{Maximum Cuts (from Coulter Beeson)}
\label{sec-4}
In this problem you'll define and show that the MaxCut problem is
NP-complete, using a reduction from a known NP-Complete problem,
Partition.  We'll first define Partition, which is already familiar to you, 
and then move on to MaxCut.

\textbf{Partition} is another name for the decision problem we saw in class,
where thieves try to evenly share their loot. An instance of this
problem is a set $I = \{v_1, v_2 ... v_n\}$ of values for $n$ items.
$I$ is a Yes-instance if the items can be fairly split, that is, if we
can partition the set $[1..n]$ into two nonempty subsets $S$ and
$\overline{S}$ such that $\sum_{i\in S}v_i = \sum_{j \in \overline{S}}
v_j$.

A \textbf{cut} in an undirected graph $G=([1..n],E)$ divides the set of nodes $[1..n]$ into
two non-empty disjoint subsets. That is, a cut is a pair $(S,
\overline{S})$ where $\overline{S} = V-S$
and $0 < |S| < n$. If weight $w(i,j)$ is associated with each edge $(i,j)$ of the graph,
then the weight of the cut $(S,\overline{S})$ is the sum of
the weights of edges that cross the cut. That is,
\[
w(S,\overline{S}) = \sum_{i \in S, j \in \overline{S}} w(i,j).
\]
We would like to show that there is there is unlikely to be an
efficient algorithm that takes as input a graph $G$ and weight
function $w$, and finds a cut of maximum weight. However
NP-Completeness is only defined for decision problems.

\begin{enumerate}
\item \question{Write down a decision version of the problem, which
we'll call MaxCut. What is an instance $I$ of MaxCut? Which are the
Yes-instances of MaxCut?}

\item \question{Show that $\text{MaxCut} \in \text{NP}$.}

\item Here you'll reason about a useful building block for a reduction
from Partition to MaxCut.
Let $I = \{v_1,v_2,\ldots,v_n\}$ be an instance of Partition, with
$\sum_{i=1}^n v_i = 2V$.  From $I$ we'll build a graph $G =
   ([1..n], E)$ where $E$ has an undirected edge between every pair of
nodes $(i,j)$, with weight $w(i,j) = v_i v_j$.
\question{Show that if $(S, \overline{S})$ is a partition of $[1..n]$ such that
$\sum_{i\in S} v_i = V$, then in $G$, the weight of the cut $(S,S')$
is $V^2$.}

\item \question{Show that if $(S, \overline{S})$ is a partition of $[1..n]$ such that
$\sum_{i\in S} v_i \neq V$, then in the graph $G$ of part 3, the weight of the cut $(S,\overline{S})$
is less than $V^2$.}
\emph{Hint:} use the fact that if $x + y = 2V$ then $x y$ attains its unique maximum
when $x = y = V$.

\item \question{Describe an efficient reduction that maps instances $I$ of Partition to instances $I'$ 
of MaxCut.}
You can describe your reduction in words or implement it in pseudocode.

\item \question{Explain why your reduction runs in polynomial time.}

\item \question{Show that if $I$ is a Yes-instance of Partition then $I'$ is a
Yes-Instance of MaxCut.} (Part 3 should be useful here, as well as
your definition of part 1.)

\item \question{Show that if $I'$ is a Yes-instance of MaxCut then $I$ is a Yes-Instance of Partition.}
(Part 4 should be useful here.)

\item In class we saw a dynamic programming algorithm for the Partition
problem, with running time $\Theta(nV)$, where $2V = \sum_{i=1}^n v_i$.
\question{Why does this not imply that $\text{P}=\text{NP}$?} (Think what is
the size of a Partition instance; is the running time of the dynamic
programming algorithm polynomial or exponential in the input size?)

\item \textbf{Bonus} \question{Prove the hint given in part 4.}
\end{enumerate}
\end{document}
